% !TEX root = ../../main.tex
\section{Phạm vi và đối tượng nghiên cứu}
\label{sec:pham_vi}
Đề tài xác định rõ đối tượng và phạm vi nghiên cứu:
\begin{itemize}
    \item \textbf{Đối tượng nghiên cứu:} Quy trình nghiên cứu khoa học, đọc hiểu tài liệu chuyên sâu, trích xuất mạng lưới tri thức học thuật và kỹ thuật biên dịch bản thảo khoa học chuẩn quốc tế.
    \item \textbf{Khách thể thụ hưởng:} Giảng viên, nghiên cứu sinh, học viên cao học và sinh viên thực hiện nghiên cứu khoa học, khóa luận và đồ án tốt nghiệp tại VinUniversity và các cơ sở giáo dục đại học định hướng nghiên cứu.
    \item \textbf{Định dạng tài liệu hỗ trợ:} Ấn phẩm học thuật chuẩn PDF (đặc biệt là bài báo 2 cột tiếng Anh và tiếng Việt), mã nguồn bản thảo khoa học \LaTeX\ và định dạng trích dẫn chuẩn BibTeX.
    \item \textbf{Hạ tầng kỹ thuật triển khai:} Máy chủ đám mây Linux (Ubuntu 24.04 LTS), hệ quản trị PostgreSQL 16 tích hợp \texttt{pgvector}, môi trường hộp cát E2B Sandbox và nền tảng giám sát LLM Langfuse.
\end{itemize}
